\subsection{An exact covering candidate: the two-step origami}
\label{sec:origami-covariance}

We now test an actual degree-four covering family, using the
full profinite completions throughout. The source model and its
monodromy words are those of
\cite[Proposition 4.2 and \S7.2]{Moeller2005}; the comparison
below is proved from the marked cyclic maps of B.1.7.
Let \(U=\mathbb P^1\setminus\{0,1,\infty\}\), let
\(P(x)=4x(x-1)\), and take the family
\[
 X_t:\quad y^2=(P(x)^2-(1-t))(P(x)-1).
                                                        \tag{B.O1}
\]
It gives the geometric monodromy homomorphism
\[
 j_*:\widehat{\pi_1(U,\overrightarrow{01})}
             \longrightarrow\widehat\Gamma_{2,0}.
                                                        \tag{B.O2}
\]
The ambiguity in a path to the target base point is simultaneous
conjugation in \(\widehat\Gamma_{2,0}\). Choose the chain
\(\alpha_1,\ldots,\alpha_5\) and the quilt of that source,
so that the two positive free generators have images
\[
 X=j_*(x)=\alpha_1^2\alpha_3\alpha_5^2,
 \qquad Y=j_*(y)=\alpha_2\alpha_4.                       \tag{B.O3}
\]
These words can also be read from the three vertical and two
horizontal cylinder cores of the four-square presentation.

For \(F\in\widehat{GT}\), let \(\alpha_0(F)\) be the
normalized action on the free source, and let \(\alpha_2(F)\)
be any representative of the outer action supplied by A.6.1.
The candidate invariant is the orbit map
\[
 \mathcal I_{S_2}(F)=
 [\alpha_2(F)\circ j_*\circ\alpha_0(F)^{-1}]
 \quad\text{in}\quad
 \operatorname{Hom}_{\rm cont}(\widehat F_2,
               \widehat\Gamma_{2,0})/
               \operatorname{Inn}(\widehat\Gamma_{2,0}).
                                                        \tag{B.O4}
\]
This target is the compact Hausdorff quotient of
\(\widehat\Gamma_{2,0}^2\) by simultaneous conjugation.
Thus (B.O4) is an actual continuous map, with its inner
ambiguity specified.

\begin{result}{B.4.3}{Proved}{The two-step origami outer invariant is constant on full GT}
For every multiplier and every \(F\in\widehat{GT}\),
\[
                \mathcal I_{S_2}(F)=[j_*].              \tag{B.O5}
\]
The degree-four cover defining this candidate has monodromy
group of order eight and deck group of order two. Its
outer monodromy invariant therefore supplies no separation
between \(G_{\mathbb Q}\) and \(\widehat{GT}\).
\end{result}
\begin{proof}
We first retain every marking in the genus-zero comparison.
Set \(a^2=1-t\). Start from the ordered four-point sphere
with points \(0,\infty,1,1-t\) in coordinate \(z\).
The replacement of \(t\) by \(1-t\) is the actual
peripheral relabelling, which commutes with the outer pure
GT action. Pull these points back under \(z=w^2\),
retaining the two branch points. The upstairs marked set is
\[
          \{0,\infty,1,-1,a,-a\}.                       \tag{B.O6}
\]
The two points \(a,-a\) may be interchanged; the four
other names are retained. Choose the sheet \(w=1\)
over \(z=1\). In the completed centralizer, the subgroup
fixing this sheet maps isomorphically to the quotient
four-point group: the two lifts of each quotient mapping
class differ by the deck involution, and exactly one
fixes the chosen sheet. This is exactly the choice
subgroup of B.1.7. Its fewer-orbit case follows there
by adding and then forgetting a third regular orbit.
Consequently the homomorphism producing (B.O6) intertwines
the two actions of the same \(F\), up to target inner
automorphism.
The witness for this comparison fixes
\(0,\infty,1,-1\) and permits only \(a,-a\)
to be exchanged: the cyclic label correction of
B.1.7 fixes these names because \(\lambda\equiv1\pmod2\),
and the pure quotient witness lifts through the chosen-sheet
subgroup. This pins the witness inside the mixed group
on which the next forgetting is defined.

Next forget \(w=0\). The remaining five-point sphere
has two named branch points \(-1,\infty\), a named
regular point \(1\), and the unordered pair \(a,-a\).
Pull back under
\[
                   w=P(x),\qquad w+1=(2x-1)^2.          \tag{B.O7}
\]
This is again the marked cyclic cover of B.1.7. Choose
the sheet \(x=(1+\sqrt2)/2\) over \(w=1\); the
choice eliminates the order-two deck kernel as before.
We use the extended choice subgroup allowing the two
regular orbit names \(a,-a\) to be exchanged. Its
quotient is exactly the five-point mixed group just
specified: those regular orbits all have trivial cover
monodromy, and choosing the sheet above \(1\) removes
the whole deck kernel.
There are now eight marked points: the fixed points
\(1/2,\infty\) and the six inverse images of
\(1,a,-a\). Forget the two fixed points. The six
remaining points are precisely the zeros of the
right side of (B.O1).

Here the intermediate five-point group allows the
two regular orbit names \(a,-a\) to be interchanged.
We justify the use of B.1.7 on this mixed group.
On its pure open subgroup, that result gives the
same parameter with the retained labels. Its label
correction fixes the chosen sheet above \(1\), and
commutes with forgetting whole orbits. The extension
to the mixed group is unique after this inner
correction: an automorphism trivial on the pure
subgroup has each \(g^{-1}\beta(g)\) centralizing it.
That centralizer is trivial. Indeed its permutation
would act trivially in the natural
\(\mathfrak S_5\)-factor of
\(\operatorname{Out}(\widehat\Gamma_{0,5})\), which
is faithful by \cite[Corollary 2.8]{CA_HMM}; after
the permutation vanishes, centrelessness of the pure
group kills the element. Thus the actual mixed
cover and both forgetful maps retain the asserted
outer comparison.
After the first forgetting its witness \(h_1\)
lies in this five-point mixed domain. If \(L_2\)
denotes the second lift followed by its forgetting,
and its witness is \(h_2\), the composite witness
is the single element \(h_2L_2(h_1)\). Hence the
comparison holds on the entire free source with
one inner automorphism.

Let
\(q:\widehat\Gamma_{2,0}\twoheadrightarrow
\widehat\Gamma_{0,[6]}\) be the hyperelliptic quotient.
The preceding two cyclic lifts and two forgetful maps
compute \(qj_*\) on the entire source group. B.4.2
identifies the action induced by \(\alpha_2(F)\)
through \(q\) with the standard six-point action of
the same \(F\). We therefore have an inner comparison
between \(q\alpha_2(F)j_*\) and
\(qj_*\alpha_0(F)\). Lift its conjugating element
through the surjection \(q\), and change the
representative \(\alpha_2(F)\) by that inner
automorphism. The two resulting homomorphisms agree
after \(q\). Its central kernel is
\(\langle h\rangle=C_2\), so their discrepancy is
an actual continuous homomorphism
\[
 \epsilon_F:\widehat F_2\longrightarrow\langle h\rangle,
 \qquad
 \alpha_2(F)(j_*(g))=
 \epsilon_F(g)\,j_*(\alpha_0(F)(g)).                    \tag{B.O8}
\]
Centrality proves multiplicativity directly by applying
the two homomorphisms to a product.

We now prove \(\epsilon_F=1\) without discarding this
kernel. Both products in (B.O3) consist of commuting
twists on disjoint curves. Compatibility with labelled
sewing in A.6.1 sends each such product to a conjugate
of its \(\lambda(F)\)-power, using one conjugator
for the whole multicurve. The normalized free action
also sends each of \(x,y\) to a conjugate of its
\(\lambda(F)\)-power. Thus, for \(g=x\) or \(y\),
the two sides of (B.O8), before multiplication by
\(\epsilon_F(g)\), are conjugate to the same
unit-power multitwist.

Use the actual finite homology homomorphism
\[
 \widehat\Gamma_{2,0}\longrightarrow
                         \operatorname{Sp}_4(\mathbb F_3).
                                                        \tag{B.O9}
\]
For a symplectic form \(J\), the twist on a vector
\(v\) acts by \(I+vv^{\mathsf T}J\). Pairwise
disjoint curves have pairings zero, so the nilpotent
parts of their transvections have all pairwise
products zero. Their multitwist, and every unit
power of it, is unipotent. Its characteristic
polynomial in (B.O9) is \((Z-1)^4\). The
hyperelliptic involution acts by \(-I\), so its
product with that multitwist has characteristic
polynomial \((Z+1)^4\). These polynomials differ
over \(\mathbb F_3\); their traces are respectively
\(1\) and \(2\). Consequently the two elements
cannot be conjugate. This forces
\(\epsilon_F(x)=\epsilon_F(y)=1\), proving (B.O5).

For completeness the covering maps of the actual
origami are
\[
\begin{aligned}
 (x,y)&\longmapsto (P(x),(2x-1)y),\\
 (w,v)&\longmapsto (w^2,wv),
\end{aligned}                                                    \tag{B.O10}
\]
with intermediate model
\(v^2=(w-1)(w+1)(w^2-(1-t))\) and final elliptic
model \(v^2=u(u-1)(u-(1-t))\). Substitution proves
both morphism identities. Its four-sheet monodromies
are \((12)(34)\) and \((23)\); they generate the
dihedral group of order eight, and their centralizer
in \(\mathfrak S_4\) has order two. The deck group
is that centralizer. A degree-four quotient by its
deck group would require four deck transformations,
which this explicit count excludes. Nevertheless,
(B.O10) retains its two actual quadratic stages.
\end{proof}

The programs \path{code/check_origami_s2.py} and
\path{code/check_origami_central_detector.py}, with
their actual outputs, check the two polynomial
identities, every permutation in the deck
centralizer, and the integral symplectic chain
matrices before reduction modulo three. The latter
checks both unit residues. Equation (B.O5) is the
outer invariant (B.O4) for this specified family.
M\"oller's more tightly pinned relation also specifies
a tangential comparison involving the Kummer
coefficient of two; recovering that exact comparison
requires keeping its path rather than passing to
the quotient in (B.O4). We do not infer that based
identity from an equality of outer classes.
